\chapter{The terminal, byte by byte}
\label{ch:terminal}

Chapter~\ref{ch:intro} described a terminal in plain words. This
chapter looks underneath. We will see exactly what bytes travel
between your program and the terminal emulator, what those bytes mean,
and why the protocol is the way it is. By the end you will be able to
read raw terminal output, write ``Hello'' in red without any library,
and trust your mental model of what \maya{} is doing one floor below
the API.

There is nothing magical here. The terminal is one of the simplest
abstract machines in computing. It has a small alphabet and a small
list of rules; we will learn both.

\section{What is a ``state machine''?}
\label{sec:term:fsm}

The chapter is titled ``the terminal as a state machine'', so it is
worth saying what that phrase means before we use it.

A \textbf{state machine} (formally, a \emph{finite-state automaton})
is a model of computation with three parts:

\begin{enumerate}[leftmargin=*]
  \item A finite set of \textbf{states} --- things the system can
    \emph{be in} right now.
  \item An \textbf{alphabet} --- the kinds of \emph{input}
    the system reacts to.
  \item A \textbf{transition table} --- for each (state, input)
    pair, the state the machine moves to next, and any output it
    produces along the way.
\end{enumerate}

The machine sits in a state, reads an input, looks up the transition,
and moves. That is the entire model. A vending machine is one: state
is ``how much money has been inserted''; alphabet is ``coin'',
``button press''; transitions either accept the coin (raising the
balance) or, if the button is pressed and balance is sufficient,
dispense a snack and return to zero.

A terminal emulator is a state machine. Its states include:

\begin{itemize}[leftmargin=*]
  \item \emph{Ground} --- the default, awaiting the next byte.
  \item \emph{Escape} --- the previous byte was \code{ESC};
    expecting a follow-up.
  \item \emph{CSI} --- inside a Control Sequence; reading parameter
    digits.
  \item \emph{OSC} --- inside an Operating System Command; reading
    a string up to a terminator.
  \item \emph{DCS} --- inside a Device Control String (rare).
\end{itemize}

Its alphabet is ``one byte''. Its transitions read like:

\begin{center}
\small
\begin{tabular}{lll}
\textbf{State} & \textbf{Input byte} & \textbf{Action / next state} \\
\hline
Ground  & printable & emit the glyph; stay in Ground \\
Ground  & \code{0x1B} (ESC) & move to Escape \\
Ground  & \code{0x0A} (LF) & move cursor down; stay in Ground \\
Escape  & \code{[}       & move to CSI \\
Escape  & other         & dispatch a one-byte escape; back to Ground \\
CSI     & digit         & accumulate into current parameter \\
CSI     & \code{;}       & start the next parameter \\
CSI     & letter         & dispatch the sequence; back to Ground \\
\end{tabular}
\end{center}

This is a tiny excerpt of what a real terminal emulator implements
(Paul Williams's well-known VT500 state diagram has a few more rows),
but it is the whole essence. Every byte the emulator receives moves
it through these states; the side effects --- drawing glyphs,
changing colours, moving the cursor --- happen on the transitions.

\maya{}'s input parser is also a state machine, but in reverse: it
reads bytes from the keyboard and assembles them into structured
\code{Key} events. Its states are ``Ground'', ``after ESC'', ``inside
CSI'', ``maybe Alt+key''. The grammar is the same; only the
direction of the byte flow is reversed.

\begin{idea}
The terminal is a small state machine that reads bytes and updates a
grid. Your program is a state machine in the other direction:
structured intent in, bytes out.
\end{idea}

\section{Three things called ``terminal''}

The word \emph{terminal} is overloaded. Pinning down which one we mean
is half the battle. There are three distinct things, often confused.

\paragraph{The hardware terminal.}
A physical device with a keyboard and a screen that connected to a
mainframe over a serial cable. Digital Equipment Corporation's
\textbf{VT100} (released 1978) is the famous one. Nobody owns one
today, but the byte protocol it spoke survives essentially unchanged
on every modern system. When we say a modern emulator ``speaks
ANSI'', we mean it speaks VT100's grandchildren.

\paragraph{The terminal emulator.}
A program on your machine --- iTerm2, kitty, Alacritty, GNOME
Terminal, Konsole, Windows Terminal, xterm --- that pretends to be a
VT100 (or a richer descendant). It owns a window on your desktop,
draws characters into it, and forwards your keystrokes to whatever
program is running ``inside'' it.

\paragraph{The pseudo-terminal, or \emph{pty}.}
A pair of file descriptors in the operating-system kernel that hooks
the emulator up to the program. We will look at it in a moment. The
pty is the part most people never hear named, and yet it is the part
doing the actual work of being a terminal.

When this book says ``\maya{} writes ANSI to the terminal'', it means:
\maya{} calls the operating-system primitive \code{write} on the file
descriptor numbered $1$ (called \code{stdout}); the bytes flow through
the pty; the emulator on the other end reads them and updates the
window.

\subsection{What is a file descriptor?}

If you have not run into the word before: in Unix, every open ``thing
you can read from or write to'' is named by a small integer called a
\textbf{file descriptor}. Three are conventional:

\begin{center}
\small
\begin{tabular}{rl}
\textbf{0} & \code{stdin}  -- bytes the program reads as input. \\
\textbf{1} & \code{stdout} -- bytes the program writes as ordinary output. \\
\textbf{2} & \code{stderr} -- bytes the program writes as error output. \\
\end{tabular}
\end{center}

When you run a program inside a terminal, all three are wired to the
pty. When you write \code{ls > out.txt}, the shell rewires the
program's file descriptor $1$ to a freshly opened file before starting
it. The program is none the wiser: it calls \code{write(1, ...)} as
always, and bytes happen to land in a file. This is the famous Unix
``everything is a file'' design, and it is what makes terminal
programs so flexible.

\section{The pseudo-terminal: a kitchen-pass-through}

A useful mental model for the pty: imagine a kitchen with a
pass-through window. The cook (your program) puts plates of bytes on
the counter; the server (the terminal emulator) picks them up and
delivers them to the diner (the window on your screen). When the
diner says something to the server, the server sets a slip of paper
(keystrokes) on the counter for the cook to pick up. The cook and the
server never see each other; they only exchange bytes via the
counter.

In reality the ``counter'' is a chunk of kernel memory --- a buffer
with two ends. The cook (your program) sees one end as a normal file:
\code{write} to send bytes into it, \code{read} to take bytes out of
it. The server (the emulator) sees the other end the same way. The
kernel shuffles bytes between the two ends, with a few useful
behaviours layered on (input modes, signal handling --- we will get to
those).

\begin{center}
\small
\begin{verbatim}
   +----------+        write(1)         +-----+      read     +----------+
   | program  | ----------------------> |     | ------------> | emulator |
   |          | <---------------------- | pty | <------------ |          |
   +----------+        read(0)          +-----+      write    +----------+
\end{verbatim}
\end{center}

This picture is enough for the rest of the chapter. We do not need to
look at the implementation of the pty itself.

\begin{idea}
The pty is a two-way byte pipe between your program and the terminal
emulator. \code{stdout} goes \emph{to} the emulator; \code{stdin}
comes \emph{from} it.
\end{idea}

\section{The byte language: ASCII plus ANSI}

Now the central question: what bytes does the program send, and what
do they do?

\subsection{Ordinary bytes: ASCII}

Most bytes you write are interpreted literally. The byte \code{0x48}
is the letter \code{H}; \code{0x69} is \code{i}; writing the two-byte
sequence \code{0x48 0x69} prints \code{Hi}. This mapping is
\textbf{ASCII}, a standard from 1963 covering the bytes from $0$ to
$127$. Every English letter, every digit, every common punctuation
mark has a fixed ASCII byte.

Characters outside ASCII (accented letters, CJK, emoji) are sent as
\textbf{UTF-8} sequences of two-to-four bytes per character. Modern
terminals all understand UTF-8; \maya{} produces UTF-8 throughout. We
will return to the question of \emph{cell width} for non-ASCII
characters in the canvas chapter.

\subsubsection*{A quick word on UTF-8}

UTF-8 is the variable-length encoding of Unicode that the world
standardised on. The rule, in one paragraph: code points 0--127 (the
old ASCII range) are encoded as a single byte, exactly as in ASCII;
code points 128--2047 take two bytes; 2048--65535 take three bytes;
the remainder (including emoji and obscure scripts) take four. The
first byte of a multi-byte sequence has a specific high-bit pattern
identifying the length (\code{110xxxxx} for two-byte, \code{1110xxxx}
for three-byte, \code{11110xxx} for four-byte), and subsequent bytes
all begin with \code{10xxxxxx}.

Example. The letter \code{H} is \code{0x48} (one byte). The Greek
letter mu, $\mu$ (U+03BC), takes two: \code{0xCE 0xBC}. The Chinese
character zh\=ong (U+4E2D) takes three: \code{0xE4 0xB8 0xAD}. A
grinning emoji (U+1F600) takes four: \code{0xF0 0x9F 0x98 0x80}.

Why this matters for a TUI: the number of \emph{bytes} in a string
is not the number of \emph{characters}, and the number of
characters is not the number of \emph{cells} the terminal will use
to display it. Three different quantities. \maya{}'s text layer
distinguishes them carefully; you usually do not have to.

\subsection{Control bytes}

A handful of bytes in the range $0$--$31$ have \emph{control}
meanings: they do not print a glyph, they tell the terminal to do
something. The interesting ones:

\begin{center}
\small
\begin{tabular}{rll}
\textbf{Byte (hex)} & \textbf{Name} & \textbf{Effect} \\
\hline
\code{0x07} & BEL & beep / bell \\
\code{0x08} & BS  & backspace --- cursor left one cell \\
\code{0x09} & HT  & horizontal tab --- advance to next tab stop \\
\code{0x0A} & LF  & line feed --- cursor down one row \\
\code{0x0D} & CR  & carriage return --- cursor to column $1$ \\
\code{0x1B} & ESC & escape --- start of an escape sequence \\
\code{0x7F} & DEL & delete --- treated as backspace on input \\
\end{tabular}
\end{center}

You have met most of these. \code{LF} is the newline. \code{CR} is
why old-school progress bars work --- print \code{CR}, then a fresh
line of text, and you overwrite the previous one in place.

The interesting byte is \code{ESC}, hex \code{0x1B}, decimal $27$. It
is the door into the ANSI escape sequence language.

\subsection{Anatomy of an escape sequence}

An ANSI escape sequence has a strict shape:

\begin{center}
\small
\begin{tabular}{l}
\code{ESC} \, \code{[} \, \emph{parameters} \, \emph{final-byte} \\
\end{tabular}
\end{center}

\begin{itemize}[leftmargin=*]
  \item \code{ESC} is the single byte \code{0x1B}.
  \item \code{[} is the literal ASCII byte \code{0x5B}. Together,
    \code{ESC [} (two bytes) form the \textbf{Control Sequence
    Introducer}, abbreviated \textbf{CSI}.
  \item \emph{parameters} is a list of decimal numbers separated by
    semicolons. Each number is written as its ASCII digits ---
    so the number $31$ is two bytes, \code{0x33 0x31}, not the binary
    value $31$.
  \item \emph{final-byte} is a single ASCII letter that says which
    command this sequence is. The letter chooses which family the
    command belongs to.
\end{itemize}

Here is a complete example, byte by byte. Suppose we want to apply
the colour ``bold red'':

\begin{terminalbox}
ESC  [   1   ;   3   1   m
0x1b 0x5b 0x31 0x3b 0x33 0x31 0x6d
\end{terminalbox}

Seven bytes. The first two are CSI. \code{1 ; 3 1} is the parameter
list ``$1$, $31$''. The final byte \code{m} says ``this is an
\textbf{SGR} command --- Select Graphic Rendition''. Parameter $1$
within SGR means ``bold''; $31$ means ``red foreground''. The terminal
emulator parses the seven bytes, updates its internal style for
\emph{future} text, and waits. The next ordinary bytes you send are
drawn bold and red.

If we now write \code{Hi}, then send \code{ESC [ 0 m} to reset, the
full transcript is:

\begin{terminalbox}
\textbackslash x1b[1;31mHi\textbackslash x1b[0m
\end{terminalbox}

This is what \maya{} ultimately produces. Every layer of the library
exists to spare you from writing those bytes by hand.

\begin{idea}
An escape sequence is \code{ESC [}, then a few numbers with
semicolons, then one letter. The letter picks the family; the numbers
parameterise it.
\end{idea}

\subsection{The families}

The final byte categorises the command. The families you will see in
\maya{}'s output, in rough order of frequency:

\begin{center}
\small
\begin{tabular}{lll}
\textbf{Family} & \textbf{Final byte} & \textbf{Purpose} \\
\hline
SGR             & \code{m} & graphic rendition: colour, bold, italic, etc. \\
CUP             & \code{H} or \code{f} & move cursor to (row, column) \\
CUU / CUD       & \code{A} / \code{B} & cursor up / down by $n$ \\
CUF / CUB       & \code{C} / \code{D} & cursor forward / back by $n$ \\
CHA             & \code{G} & cursor to absolute column \\
ED              & \code{J} & erase part of the display \\
EL              & \code{K} & erase part of the line \\
DECSET / DECRST & \code{h} / \code{l} & set / reset private modes (with \code{?}) \\
\end{tabular}
\end{center}

Three additional things to know.

\paragraph{ANSI columns and rows are 1-based.}
\code{ESC [ 5 ; 12 H} moves the cursor to row $5$, column $12$, where
the top-left cell of the screen is $(1, 1)$, \emph{not} $(0, 0)$. This
is a constant source of off-by-one bugs; \maya{} converts to/from
$0$-based internally and \code{move\_to}'s signature is documented as
1-based at the boundary.

\paragraph{Private modes use \code{?}.}
DEC's extensions to ANSI use a question mark inside the parameter list
to distinguish ``private'' parameters that only DEC-style terminals
understand. \code{ESC [ ? 25 h} (show cursor) and
\code{ESC [ ? 1049 h} (enter alternate screen) are private; the plain
SGR \code{ESC [ 1 m} is not.

\paragraph{Numbers default to $1$.}
If you write \code{ESC [ A} (just the family with no parameter), the
parameter is implicitly $1$ --- ``cursor up one row''. Most families
follow this convention; \maya{} explicitly emits the number anyway,
because being explicit avoids surprises across emulators.

\subsection{SGR: colours and attributes}

SGR is the family you will see most. It composes: \code{ESC [ 1 ; 31
; 4 m} applies bold, red, and underline together. Then ordinary text
is drawn with all three until another SGR command changes them.

The parameters you will care about:

\begin{lstlisting}
ESC [ 0 m              // reset every attribute
ESC [ 1 m              // bold
ESC [ 2 m              // dim
ESC [ 3 m              // italic
ESC [ 4 m              // underline
ESC [ 7 m              // inverse (swap fg / bg)
ESC [ 9 m              // strikethrough

ESC [ 22 m             // bold + dim off
ESC [ 23 m             // italic off
ESC [ 24 m             // underline off
ESC [ 27 m             // inverse off
ESC [ 29 m             // strikethrough off

ESC [ 30..37 m         // 3-bit foreground (black..white)
ESC [ 40..47 m         // 3-bit background
ESC [ 90..97 m         // bright fg
ESC [ 100..107 m       // bright bg
ESC [ 39 m             // default fg
ESC [ 49 m             // default bg

ESC [ 38 ; 5 ; N m     // 256-colour fg, palette index N
ESC [ 38 ; 2 ; R ; G ; B m   // 24-bit truecolour fg
ESC [ 48 ; 2 ; R ; G ; B m   // 24-bit truecolour bg
\end{lstlisting}

Modern terminals all support truecolour (24 bits = 16,777,216 colours).
\maya{}'s \code{Fg<R, G, B>} compiles to exactly \code{ESC [ 38 ; 2 ;
R ; G ; B m}. If you read \file{maya/include/maya/terminal/ansi.hpp}
you will find these as \code{constexpr std::string\_view} constants:

\begin{lstlisting}
inline constexpr std::string_view bold      = "\x1b[1m";
inline constexpr std::string_view italic    = "\x1b[3m";
inline constexpr std::string_view underline = "\x1b[4m";
inline constexpr std::string_view reset     = "\x1b[0m";
\end{lstlisting}

\code{constexpr} means the string literally lives in the binary; no
allocation, no construction at startup. We will explain
\code{string\_view} and \code{constexpr} properly in
Chapter~\ref{ch:cpp}.

\subsection{Cursor movement and erasing}

A short worked example: print ``A'' at row $3$, column $5$, then
``B'' at row $3$, column $20$:

\begin{terminalbox}
ESC [ 3 ; 5 H A ESC [ 3 ; 20 H B
\end{terminalbox}

Two CUP sequences sandwich the printable letters. Each is six bytes
(seven, with the two-digit column): the cost of cursor positioning is
exactly the length of the decimal numbers plus four (\code{ESC},
\code{[}, \code{;}, \code{H}). This is why diff-based rendering is so
effective: most cursor moves are tiny.

For erasing, the three useful ED forms:

\begin{lstlisting}
ESC [ 0 J     // erase from cursor to end of screen
ESC [ 1 J     // erase from start of screen to cursor
ESC [ 2 J     // erase entire screen
\end{lstlisting}

And EL, line-erasing:

\begin{lstlisting}
ESC [ 0 K     // erase from cursor to end of line
ESC [ 1 K     // erase from start of line to cursor
ESC [ 2 K     // erase entire line
\end{lstlisting}

The cursor itself does not move when you erase --- the cells under the
cursor change to blank but the position stays. If you want the
classic \emph{clear-screen-and-home} effect, you send both:
\code{ESC [ 2 J ESC [ H} (the \code{H} with no parameters means ``move
to row $1$, column $1$'').

\begin{warning}
Erasing does not reset the colour. A cell ``erased'' to blank keeps
its background colour. If you erased after setting a red background,
you get a red rectangle. To truly blank the screen to the default,
send \code{ESC [ 0 m} (reset colours) \emph{before} erasing.
\end{warning}

\subsection{Modes: cursor, alt-screen, mouse, paste}

A different style of command toggles \emph{modes}: terminal-wide flags
the emulator carries between sequences. These use the DEC private form
with a \code{?}:

\begin{lstlisting}
ESC [ ? 25 l           // hide the cursor
ESC [ ? 25 h           // show the cursor
ESC [ ? 1049 h         // enter alternate screen
ESC [ ? 1049 l         // leave alternate screen
ESC [ ? 1000 h         // mouse press/release events
ESC [ ? 1002 h         // mouse + drag motion events
ESC [ ? 1006 h         // SGR-encoded mouse coordinates (no 223-col limit)
ESC [ ? 2004 h         // bracketed paste (paste arrives wrapped)
ESC [ ? 2026 h         // begin synchronised update
ESC [ ? 2026 l         // end synchronised update
\end{lstlisting}

A few of these deserve names of their own.

\paragraph{The alternate screen.}
When you run \code{vim} or \code{less}, you notice that quitting
restores whatever you had on screen \emph{before} the editor started
--- the editor's contents do not pollute your scrollback. This is the
alternate screen. The emulator keeps two screen buffers; \code{?1049
h} switches to the secondary one (saving the primary), \code{?1049 l}
switches back. \maya{} programs in fullscreen mode emit the enter
sequence at startup and the leave sequence at exit, which is why they
behave the same way.

\paragraph{Bracketed paste.}
When the user pastes a chunk of text, the emulator wraps it with the
sentinels \code{ESC [ 200 \~{}} (before) and \code{ESC [ 201 \~{}}
(after). Without this, a pasted multi-line block looks identical to
the user typing very fast, and a TUI editor cannot tell the
difference. With it, the editor sees: ``a paste begins, here are the
bytes, the paste ends'' --- and can treat the whole block atomically
(no auto-indent per line, for example). \maya{} surfaces this as a
\code{Paste} variant of its \code{Event} type.

\paragraph{Synchronised update.}
A new and very useful mode. When \code{?2026 h} is sent, the terminal
buffers all output until \code{?2026 l}. The frame is then drawn
atomically --- no tearing, no half-painted UI even on a slow connection.
Modern terminals (kitty, WezTerm, foot, Ghostty, recent xterm)
support this; older ones ignore the sequence harmlessly. \maya{}'s
\code{env\_supports\_synchronized\_output} in
\file{terminal/ansi.hpp} detects this from environment variables and
decides whether to wrap each frame in sync brackets.

\subsection{Encoding non-ASCII keys: the Kitty Keyboard Protocol}

Plain ANSI cannot tell \code{Ctrl+M} from \code{Enter} (both arrive as
byte $0x0D$), or \code{Shift+Tab} from a regular escape sequence. This
has been a chronic source of bugs in TUI editors for thirty years.
The \textbf{Kitty Keyboard Protocol} (KKP) and the older
\code{modifyOtherKeys=2} extension fix this by sending fully
disambiguated key reports.

\maya{} opts in on startup with two sequences:

\begin{lstlisting}
ESC [ > 1 u            // push KKP flag 1 ("disambiguate")
ESC [ > 4 ; 2 m        // xterm modifyOtherKeys=2
\end{lstlisting}

and opts out on exit with \code{ESC [ < u} and \code{ESC [ > 4 m}.
Terminals that do not understand the sequences ignore them silently;
the worst case is the older behaviour. We discuss the input side
of this in Chapter~\ref{ch:events}.

\section{Cooked mode, raw mode, and \code{termios}}

Now we turn from output to input. In a normal shell, what happens
when you type a line of text?

The kernel sits between your keyboard and any program reading
\code{stdin}. By default it is in \textbf{canonical} (cooked) mode:

\begin{itemize}[leftmargin=*]
  \item Keystrokes are buffered until you press \code{Enter}, then
    delivered as a whole line.
  \item Backspace edits the buffer in place (the program never sees
    the deletion).
  \item \code{Ctrl-C} is intercepted and sent as a \emph{signal}
    (\code{SIGINT}) rather than as data.
  \item \code{Ctrl-D} on an empty line is interpreted as
    end-of-input.
  \item Each typed character is also \emph{echoed} back to the
    screen automatically.
\end{itemize}

This is exactly the behaviour you want at a shell prompt. It is
exactly the wrong behaviour for a TUI. A TUI wants every keystroke
delivered immediately, byte for byte, with no automatic echo and no
special handling of control characters.

The opposite mode is \textbf{raw mode}, configured via the
\code{termios} structure --- a C \code{struct} the kernel exposes
through \code{tcgetattr} and \code{tcsetattr}.

\subsection{The \code{termios} struct}

\code{termios} has four flag fields plus a small array of special
characters. Each flag field is a bitmask: bits enabled mean different
behaviours active.

\begin{center}
\small
\begin{tabular}{ll}
\textbf{Field} & \textbf{Meaning} \\
\hline
\code{c\_iflag} & input-processing flags \\
\code{c\_oflag} & output-processing flags \\
\code{c\_cflag} & control flags (baud rate, parity --- mostly serial relics) \\
\code{c\_lflag} & local flags (canonical mode, echo, signals) \\
\code{c\_cc[]}  & special characters (which byte is INTR? EOF? \dots) \\
\end{tabular}
\end{center}

To enter raw mode, you turn off a specific list of bits:

\begin{lstlisting}[language=C]
#include <termios.h>
#include <unistd.h>

termios cooked, raw;
tcgetattr(STDIN_FILENO, &cooked);   // save the current settings
raw = cooked;                       // copy, mutate the copy

raw.c_lflag &= ~(ICANON | ECHO | ISIG | IEXTEN);
raw.c_iflag &= ~(IXON | ICRNL | INPCK | ISTRIP);
raw.c_oflag &= ~(OPOST);
raw.c_cc[VMIN]  = 0;     // non-blocking reads
raw.c_cc[VTIME] = 0;
tcsetattr(STDIN_FILENO, TCSAFLUSH, &raw);

// ... TUI runs here ...

tcsetattr(STDIN_FILENO, TCSAFLUSH, &cooked);   // restore
\end{lstlisting}

The bits, decoded:

\begin{itemize}[leftmargin=*]
  \item \code{ICANON} --- canonical mode (line buffering). Off means
    every keystroke comes through immediately.
  \item \code{ECHO} --- echo input. Off means typed characters do not
    appear on screen automatically; the program decides.
  \item \code{ISIG} --- signal generation on \code{Ctrl-C} /
    \code{Ctrl-Z}. Off means those become normal bytes ($0x03$,
    $0x1A$).
  \item \code{IEXTEN} --- implementation-defined extensions
    (\code{Ctrl-V} as ``literal next'', etc.). Off, just to be
    safe.
  \item \code{IXON} --- flow control (\code{Ctrl-S} freezes the
    terminal, \code{Ctrl-Q} unfreezes). Off; you want those bytes
    yourself.
  \item \code{ICRNL} --- translate \code{CR} ($0x0D$) into \code{LF}
    ($0x0A$) on input. Off, so you see exactly what was typed.
  \item \code{INPCK}, \code{ISTRIP} --- parity checking and 7-bit
    stripping. Both relics from serial-line days.
  \item \code{OPOST} --- output processing, mostly translating
    \code{LF} into \code{CR LF}. Off, because you will be writing
    full escape sequences yourself.
  \item \code{VMIN} = 0, \code{VTIME} = 0 --- ``return from
    \code{read} immediately if there is no input''. The TUI loop
    polls; it does not block.
\end{itemize}

\begin{warning}
You must restore \code{cooked} on exit. If the program crashes before
restoring, the user's shell will appear broken --- typed characters
will not echo, \code{Enter} will not commit, \code{Ctrl-C} will not
work. The remedy is the magic command \code{stty sane}, which resets
\code{termios} to sensible defaults. \maya{} uses a RAII guard (see
Chapter~\ref{ch:cpp}) so the destructor restores cooked mode even on
exceptions.
\end{warning}

\subsection{What keystrokes look like in raw mode}

In raw mode, a single keypress arrives as one or more bytes.

\begin{itemize}[leftmargin=*]
  \item Printable keys: one byte, their ASCII value. \code{a} arrives
    as $0x61$.
  \item Control keys: \code{Ctrl-A} is $0x01$, \code{Ctrl-B} is
    $0x02$, \dots, \code{Ctrl-Z} is $0x1A$. (The pattern: clear the
    top three bits of the letter.)
  \item \code{Enter}: $0x0D$ ($CR$).
  \item \code{Tab}: $0x09$ ($HT$).
  \item \code{Backspace}: $0x7F$ ($DEL$).
  \item \code{Escape}: $0x1B$. By itself \emph{or} as the start of a
    sequence.
  \item Function and arrow keys: an escape sequence. \code{Up}
    arrives as the three bytes \code{ESC [ A}; \code{Down} as
    \code{ESC [ B}; \code{Right} \code{ESC [ C}; \code{Left}
    \code{ESC [ D}.
  \item \code{Alt-}\emph{x}: the bytes \code{ESC} then \emph{x}.
\end{itemize}

The clash is obvious: \code{ESC} alone is a real keypress (some
people bind it), but \code{ESC} is \emph{also} the first byte of every
arrow-key sequence. The standard trick is to read \code{ESC},
\emph{wait a few milliseconds}, and if no more bytes arrive, treat it
as a bare \code{Escape}. If they do, parse the sequence. \maya{}'s
input parser (\file{src/terminal/input.cpp}) implements exactly this;
the timeout is around 50 ms by default.

\section{The screen as a grid}

We have looked at bytes; now look at what those bytes update.

The emulator maintains, in its own memory, a 2-D grid of \textbf{cells}.
Each cell holds:

\begin{itemize}[leftmargin=*]
  \item One Unicode \emph{grapheme} (for most languages, one code
    point --- a letter or a digit; for emoji and some scripts, a
    short cluster of code points that displays as one glyph).
  \item A foreground colour and a background colour.
  \item A set of attribute bits: bold, italic, underline, inverse,
    strikethrough, dim.
\end{itemize}

The cursor has a position (row, column) and a visibility flag. A
\textbf{frame} of UI is a complete specification of every cell that
should be on screen.

When the user resizes the window, the emulator sends the program a
notification: on Linux/macOS, the \code{SIGWINCH} signal; on Windows,
a console-resize event. The program re-queries the new size with
\code{ioctl(TIOCGWINSZ)} and re-lays-out its UI. \maya{} does this
inside the runtime; you do not need to call anything.

\subsection{Wide characters}

A subtle wrinkle: not every Unicode character takes \emph{one} cell.

\begin{itemize}[leftmargin=*]
  \item CJK ideographs (U+4E2D \emph{zh\=ong}, U+3042 \emph{hiragana
    a}) take \textbf{two} cells. Their glyph is roughly square; in a
    monospace font each one is twice as wide as a Latin letter.
  \item Many emoji also take two cells (U+1F600 smiley face).
  \item Combining marks --- a \code{COMBINING ACUTE ACCENT} on top
    of an \code{e} to render \emph{\'e} --- take \textbf{zero}
    cells; they decorate the previous cell.
\end{itemize}

If the program writes a wide character but counts it as one cell, the
emulator and the program disagree about the cursor position after,
and the layout drifts. The Unicode Consortium publishes
\file{EastAsianWidth.txt} listing the width of every code point;
\maya{} ships a generated table at
\file{include/maya/text/unicode\_width\_table.hpp} and consults it
when laying out text. You do not call into the table directly ---
\maya{}'s text widget does it for you.

\section{Signals: \code{SIGWINCH}, \code{SIGINT}, and friends}
\label{sec:term:signals}

A \textbf{signal} on Unix is a tiny asynchronous notification the
kernel delivers to a process. The process can install a handler ---
a function that runs when the signal arrives --- or it can use the
default behaviour (often: ``terminate''). Two signals matter to a
TUI.

\paragraph{\code{SIGWINCH} --- window-size change.}
When the user resizes the terminal window, the kernel delivers
\code{SIGWINCH} to whichever process owns the foreground of the pty.
The signal carries no payload --- the process must re-ask the kernel
for the new size. The mechanism is the system call \code{ioctl},
broadly ``do a non-standard thing on a file descriptor''. The
request \code{TIOCGWINSZ} (``terminal IO control: get window size'')
fills a small struct with row and column counts:

\begin{lstlisting}[language=C]
#include <sys/ioctl.h>
struct winsize ws;
ioctl(STDOUT_FILENO, TIOCGWINSZ, &ws);
// ws.ws_row, ws.ws_col, ws.ws_xpixel, ws.ws_ypixel
\end{lstlisting}

\maya{}'s runtime installs a \code{SIGWINCH} handler that simply
sets an atomic flag; the main loop checks the flag, re-queries the
size, and re-runs layout. Doing real work inside a signal handler
is dangerous (most library functions are not safe to call from one);
flagging-then-handling is the standard idiom.

\paragraph{\code{SIGINT} --- interrupt.}
Delivered by the kernel when the user presses \code{Ctrl-C} \emph{in
cooked mode}. In raw mode (which \maya{} sets), \code{Ctrl-C}
becomes the byte \code{0x03}, which the program sees as ordinary
input and may interpret however it likes. \maya{} also installs a
backup \code{SIGINT} handler so that if something forces cooked
mode back on mid-run (a misbehaving sub-process, for example), a
Ctrl-C still cleans up the terminal state before exiting.

\paragraph{\code{SIGTERM}, \code{SIGHUP}, \code{SIGSEGV}.}
\maya{}'s terminal driver installs handlers for all three so that
catastrophic exits still restore the cooked-mode terminal
before the process dies. Without this, killing a TUI with
\code{kill <pid>} or letting it crash would leave the user's shell
in raw mode. RAII is not enough alone; the destructor only runs if
the stack unwinds, and \code{SIGSEGV} does not unwind. The signal
handler is a belt-and-braces backup.

\subsection{What \code{ioctl} actually is}

The ``standard'' file operations are \code{open}, \code{read},
\code{write}, \code{close}. \code{ioctl} (``IO control'') is the
escape hatch for everything else: ``perform request \code{X} on file
descriptor \code{fd}, with these arguments''. Hundreds of requests
exist; each one is essentially a separate API hiding behind a
common name. You meet \code{ioctl} mostly for terminals
(\code{TIOCGWINSZ}), sockets, and device drivers.

In portable code you would prefer dedicated APIs; in practice no
such dedicated API exists for ``get terminal size'', so \code{ioctl}
with \code{TIOCGWINSZ} is the only path.

\section{Mouse reporting, on the wire}
\label{sec:term:mouse}

Mouse input is also delivered as escape sequences. When mouse
reporting is enabled (\code{ESC [ ?1000 h} or one of its richer
variants), every mouse event arrives as a CSI sequence the program
must parse.

The modern format is \textbf{SGR mouse reporting} (\code{?1006 h}):

\begin{center}
\small
\code{ESC [ < button ; column ; row M}~~(press / drag) \\
\code{ESC [ < button ; column ; row m}~~(release)
\end{center}

\code{button} is a small bitfield: low bits encode left/middle/right
(0/1/2); higher bits flag \code{Shift}, \code{Alt}, \code{Ctrl}
modifiers (4, 8, 16); a bit at position 32 marks ``motion event''
(this is a drag, not a click); a bit at position 64 marks wheel
scroll. \code{column} and \code{row} are 1-based.

Example. A left-click at (col 12, row 5):

\begin{terminalbox}
ESC [ < 0 ; 12 ; 5 M
\end{terminalbox}

A Shift-drag of the same button to (15, 5):

\begin{terminalbox}
ESC [ < 36 ; 15 ; 5 M
\end{terminalbox}

(Button code $36 = 0 + 4 \text{ (Shift)} + 32 \text{ (motion)}$.)

A wheel-scroll down at (40, 10):

\begin{terminalbox}
ESC [ < 65 ; 40 ; 10 M
\end{terminalbox}

(Button code $65 = 1 + 64 \text{ (wheel)}$; the low bit on wheel
means ``down''. Wheel up would be 64; left/right wheel use 66/67.)

\maya{}'s event parser decodes these into structured
\code{Mouse} events. The widget layer then routes them by position
to the widget under the cursor.

\subsubsection*{Why the SGR form?}

The old X10 mouse-reporting format (\code{?1000 h}) encoded the
column and row as \emph{single bytes}, with an offset of 32 (so the
byte for column 1 was 33, the byte for column 12 was 44). It broke
on column or row numbers above $223 - 32 = 191$, because the bytes
were only 7-bit-safe and the offsetting overflowed. Modern wide
terminals routinely exceed 200 columns; the SGR form, with decimal
coordinates, has no such limit. \maya{} enables both \code{?1000 h}
and \code{?1006 h} on startup, the latter overrides; emulators that
do not understand \code{?1006 h} fall back gracefully to the older
encoding.

\section{Bracketed paste, on the wire}
\label{sec:term:paste}

When \code{?2004 h} is enabled and the user pastes a block of text
into the terminal, the emulator wraps the pasted bytes in two
sentinel sequences:

\begin{terminalbox}
ESC [ 200 \textasciitilde~ ... pasted bytes ... ESC [ 201 \textasciitilde
\end{terminalbox}

The content between can be arbitrary, including newlines and other
control characters --- the program should treat the whole block
as one atomic event.

Example. Pasting the two lines

\begin{terminalbox}
foo \\
bar
\end{terminalbox}

produces on the wire:

\begin{terminalbox}
ESC [ 200 \textasciitilde~f o o LF b a r ESC [ 201 \textasciitilde
\end{terminalbox}

Without bracketed paste, that arrival is indistinguishable from a
very fast user typing \code{f}, \code{o}, \code{o}, \code{Enter},
\code{b}, \code{a}, \code{r}. With it, an editor knows to disable
auto-indent, treat the whole block as a single undo unit, and so
on. \maya{}'s input parser produces a single \code{Paste} event
carrying the entire content as a \code{std::string}.

\section{The input parser as a state machine}
\label{sec:term:input-fsm}

We sketched the input parser earlier as ``a state machine in the
other direction''. Concretely: it reads bytes from \code{stdin} and
builds structured \code{Event} values. Here is its core loop in
pseudo-code:

\begin{lstlisting}
enum State { Ground, AfterEsc, InCsi };
State s = Ground;
std::string params;          // accumulator for CSI parameter bytes
for (;;) {
    byte b = read_one();
    switch (s) {
      case Ground:
        if (b == 0x1B) { s = AfterEsc; start_esc_timeout(50ms); }
        else            emit_key(b);
        break;
      case AfterEsc:
        if (b == '[')  { s = InCsi; params.clear(); }
        else if (timed_out) { emit_key(Escape); s = Ground; pushback(b); }
        else            { emit_key_alt(b); s = Ground; }
        break;
      case InCsi:
        if (is_param_byte(b)) params.push_back(b);
        else if (is_final_byte(b)) {
            dispatch_csi(params, b);   // -> emit Arrow, Mouse, Paste, ...
            s = Ground;
        }
        break;
    }
}
\end{lstlisting}

A single keystroke can therefore become a single \code{Event}
(\code{A} arrives, \code{KeyKind::Char 'A'} emitted) or take a few
bytes (\code{ESC [ A} becomes \code{KeyKind::Up}) or many
(a paste of a hundred bytes is one \code{Paste} event). The 50 ms
timeout on \code{ESC} is the only piece that is not a pure function
of the byte stream --- it disambiguates a bare \code{Escape} from
an incomplete escape sequence.

\maya{}'s real parser is more elaborate (it handles modified
arrows, the Kitty Keyboard Protocol's CSI-u format,
\code{modifyOtherKeys=2}, function keys, focus events, OSC
responses), but every additional case follows the same state-machine
shape.

\section{Putting it together: a no-library Hello world}

We now know enough to write a complete TUI ``Hello, red world'' with
no library at all. It is a useful exercise: you will see exactly what
\maya{} is doing for you.

\begin{lstlisting}[language=C]
#include <stdio.h>
#include <termios.h>
#include <unistd.h>

int main(void) {
    termios cooked, raw;
    tcgetattr(STDIN_FILENO, &cooked);
    raw = cooked;
    raw.c_lflag &= ~(ICANON | ECHO);
    tcsetattr(STDIN_FILENO, TCSAFLUSH, &raw);

    // 1. Hide the cursor
    // 2. Clear the screen
    // 3. Move to row 1, column 1
    // 4. Set bold red foreground
    // 5. Write "Hi"
    // 6. Reset attributes
    // 7. Show the cursor again
    write(STDOUT_FILENO,
          "\x1b[?25l"          // hide cursor
          "\x1b[2J"            // erase entire screen
          "\x1b[1;1H"          // cursor to (row 1, col 1)
          "\x1b[1;31m"         // bold red foreground
          "Hi"
          "\x1b[0m"            // reset attributes
          "\x1b[?25h",         // show cursor
          /*len=*/ 33);

    char c;
    read(STDIN_FILENO, &c, 1);   // wait for any keypress

    tcsetattr(STDIN_FILENO, TCSAFLUSH, &cooked);
    return 0;
}
\end{lstlisting}

Save it as \file{hello.c}, build with \code{cc hello.c -o hello}, run
\code{./hello}. The screen clears; ``Hi'' appears in bold red at the
top-left; press any key; the screen returns. You have just written a
TUI by hand.

This is also a faithful caricature of what \maya{} does \emph{every
frame}. The whole library is in service of generating exactly this
style of byte stream, only longer, more carefully diffed, and
described to you as a tree of styled elements instead of a literal
string.

\section{A real maya frame, traced byte by byte}
\label{sec:term:trace}

To close the chapter, here is a hand-traced version of one
incremental redraw in a maya program. The screen currently shows
the word \code{Count: 7} on row 5, starting at column 3, in the
default white foreground. The user pressed \code{+}; maya wants the
screen to now read \code{Count: 8}. Only one cell changed --- column
10 of row 5 --- so maya emits only what is needed to update it.

The bytes \maya{} writes:

\begin{terminalbox}
\textbackslash x1b[?2026h   ~ESC [ ?2026 h ~  begin synchronised update \\
\textbackslash x1b[5;10H   ~ESC [ 5;10 H  ~  move cursor to row 5, col 10 \\
\textbackslash x1b[39;49m  ~ESC [ 39;49 m ~  default fg \& bg (style reset) \\
8                       ~~~~~~~~~~~~~~~~~~  the new digit \\
\textbackslash x1b[?2026l  ~ESC [ ?2026 l ~  end synchronised update
\end{terminalbox}

Byte-by-byte that is:

\begin{terminalbox}
1B 5B 3F 32 30 32 36 68 1B 5B 35 3B 31 30 48 1B \\
5B 33 39 3B 34 39 6D 38 1B 5B 3F 32 30 32 36 6C
\end{terminalbox}

32 bytes in total to redraw the change. Compare to what a naive
renderer would do: clear the screen, redraw every cell. For an 80
by 24 terminal that is roughly $80 \times 24 \times 10$ bytes (most
cells take around ten bytes including SGR and the character), or
close to 20\,kB. Two orders of magnitude saved by the diff.

The sync wrapper (\code{?2026 h} \dots \code{?2026 l}) is optional
and omitted by terminals that do not advertise it; \maya{}'s
renderer checks an environment-derived hint
(\code{env\_supports\_synchronized\_output} in
\file{terminal/ansi.hpp}) and elides the wrapper if the terminal
will not honour it.

The cursor positioning (\code{ESC [ 5;10 H}) and the style transition
(\code{ESC [ 39;49 m}) are produced by \maya{}'s
\code{StyleApplier::transition\_to} helper, which writes the minimal
SGR change from the previous cell's style to the new one (in this
case, the previous cell on the row was already styled the same, so
the SGR transition would actually be empty --- we include it here
for illustration).

\begin{idea}
The whole \maya{} rendering pipeline exists to make this incremental
byte stream short and correct. Everything from the canvas to the
SIMD diff to the ANSI builder is in service of: ``write the minimum
bytes needed to update only the cells that actually changed''.
\end{idea}

\begin{recap}
A terminal is a state machine on the far end of a pty. Bytes flowing
into it are mostly literal text, with a small number of control
characters and the ANSI escape language (\code{ESC [} \emph{params}
\emph{letter}) for everything else. Cooked mode lets the kernel
pre-process input; raw mode hands bytes through unchanged. The screen
is a 2-D grid of styled Unicode cells. Every layer above this in
\maya{} speaks in \cpp{} values; this one alone speaks in bytes.
\end{recap}

\section*{Workshop}

\begin{enumerate}[leftmargin=*]
  \item \textbf{Inspect your terminal's flags.} Run \code{stty -a}
    in a normal shell. Find the lines beginning with \code{iflag}
    and \code{lflag}. Identify \code{icanon} and \code{echo} (with
    or without a leading \code{-}, which means ``negated''). Those
    are the bits a TUI clears.

  \item \textbf{Identify maya's constants.} Open
    \file{maya/include/maya/terminal/ansi.hpp} (we read it above).
    Find:
    \begin{itemize}[leftmargin=*]
      \item the constant for entering the alternate screen;
      \item the constant for enabling SGR mouse reporting (1006);
      \item the constant for the Kitty Keyboard Protocol push.
    \end{itemize}
    Cross-check each one against the table earlier in this chapter.

  \item \textbf{Drive a terminal by hand.} In a scratch \cpp{} file:
    \begin{lstlisting}
#include <cstdio>
#include <unistd.h>
int main() {
    std::fputs("\x1b[?1049h", stdout);     // enter alt-screen
    std::fputs("\x1b[1;1HHello", stdout);  // print at top-left
    sleep(2);
    std::fputs("\x1b[?1049l", stdout);     // leave alt-screen
}
    \end{lstlisting}
    Build and run. Observe that your shell is exactly as you left it
    afterwards.

  \item \textbf{Type the no-library example.} Compile and run the C
    program in this chapter. Then mutate it: change \code{31}
    (red) to \code{32} (green), \code{1} (bold) to \code{4}
    (underline). Each change is one digit.

  \item \textbf{Provoke a raw-mode mistake.} Comment out the final
    \code{tcsetattr} that restores cooked mode. Recompile and run.
    Your shell will behave strangely after. Recover by typing
    \code{stty sane}, then \code{Enter} (you may need to type
    blind). You now appreciate the RAII guard.
\end{enumerate}
